\section{Desarrollo}
\subsection{Construcción de la caja}
\loremipsum

\subsection{Programación de la base de datos}
Para realizar la base de datos, se realizó un ejercicio de tablas y relaciones entre ellas (MOSTRAR ANEXO). A partir de esto, se realizarón correcciones triviales, dejando el siguiente modelo entidad-relación que construye la base de datos final \ref{fig:bd-final}.

% Simplified ER at full text width; the complete diagram (with attributes) is an annex
\end{multicols}
\begin{center}
\resizebox{\linewidth}{!}{%
\begin{tikzpicture}[
  font=\sffamily\scriptsize,
  ent/.style={rectangle, rounded corners=1.5pt, draw=lvlSectorBorder!60, fill=lvlSectorFill, line width=0.7pt, text=lvlSectorText, font=\sffamily\bfseries\small, align=center, minimum width=2.3cm, minimum height=0.85cm},
  rel/.style={diamond, aspect=2.2, draw=lvlRoleBorder!60, fill=lvlRoleFill, line width=0.6pt, text=lvlRoleText, font=\sffamily\bfseries\scriptsize, align=center, inner sep=1.2pt},
  card/.style={circle, fill=white, draw=lvlRoleBorder!40, line width=0.4pt, inner sep=1pt, text=lvlRoleText, font=\sffamily\bfseries\tiny},
  lin/.style={draw=lvlRoleBorder!55, line width=0.6pt}
]
  \node[ent] (ConsultasIA) at (0.25,-1.20) {ConsultasIA};
  \node[ent] (Usuario) at (2.30,-4.00) {Usuario};
  \node[ent] (RFID) at (0.49,-7.60) {RFID};
  \node[ent] (Equipo) at (5.08,-0.60) {Equipo};
  \node[ent] (Historial) at (4.92,-6.60) {Historial};
  \node[ent] (Caja) at (7.38,-3.20) {Caja};
  \node[ent] (Accesos) at (8.04,-9.80) {Accesos};
  \node[ent] (Controladores) at (12.30,-5.20) {Controladores};
  \node[ent] (Sensores) at (13.12,-1.40) {Sensores};
  \node[ent] (LecturasSensores) at (17.55,-0.20) {Lecturas\\Sensores};
  \node[ent] (EventosInventario) at (17.06,-4.40) {Eventos\\Inventario};
  \node[ent] (Actuadores) at (12.63,-8.60) {Actuadores};
  \node[ent] (EventosActuadores) at (16.24,-10.60) {Eventos\\Actuadores};
  \node[ent] (EtiquetasHerramienta) at (18.53,-7.40) {Etiquetas\\Herramienta};
  \node[ent] (Objeto) at (19.19,-10.60) {Objeto};
  \node[rel] (rUCo) at (1.27,-2.60) {Realiza};
  \node[rel] (rUE) at (3.69,-2.30) {Asocia};
  \node[rel] (rEC) at (6.23,-1.90) {Pertenece};
  \node[rel] (rUC) at (4.84,-3.60) {Tiene\\permiso};
  \node[rel] (rUH) at (3.61,-5.30) {Relaciona};
  \node[rel] (rHC) at (6.15,-4.90) {Relaciona};
  \node[rel] (rUR) at (1.39,-5.80) {Posee};
  \node[rel] (rRA) at (4.26,-8.70) {Usa};
  \node[rel] (rCA) at (7.71,-6.50) {Recibe};
  \node[rel] (rCCt) at (9.84,-4.20) {Administra};
  \node[rel] (rCtS) at (12.71,-3.30) {Conecta};
  \node[rel] (rCtA) at (12.46,-6.90) {Conecta};
  \node[rel] (rSL) at (15.33,-0.80) {Genera};
  \node[rel] (rSE) at (15.09,-2.90) {Detecta};
  \node[rel] (rAE) at (14.43,-9.60) {Ejecuta};
  \node[rel] (rEtE) at (17.79,-5.90) {Corresponde};
  \node[rel] (rOE) at (18.86,-9.00) {Identifica};
  \begin{scope}[on background layer]
  \draw[lin] (rUCo) -- (Usuario);
  \draw[lin] (rUCo) -- (ConsultasIA);
  \draw[lin] (rUE) -- (Usuario);
  \draw[lin] (rUE) -- (Equipo);
  \draw[lin] (rEC) -- (Equipo);
  \draw[lin] (rEC) -- (Caja);
  \draw[lin] (rUC) -- (Usuario);
  \draw[lin] (rUC) -- (Caja);
  \draw[lin] (rUH) -- (Usuario);
  \draw[lin] (rUH) -- (Historial);
  \draw[lin] (rHC) -- (Caja);
  \draw[lin] (rHC) -- (Historial);
  \draw[lin] (rUR) -- (Usuario);
  \draw[lin] (rUR) -- (RFID);
  \draw[lin] (rRA) -- (RFID);
  \draw[lin] (rRA) -- (Accesos);
  \draw[lin] (rCA) -- (Caja);
  \draw[lin] (rCA) -- (Accesos);
  \draw[lin] (rCCt) -- (Controladores);
  \draw[lin] (rCCt) -- (Caja);
  \draw[lin] (rCtS) -- (Controladores);
  \draw[lin] (rCtS) -- (Sensores);
  \draw[lin] (rCtA) -- (Controladores);
  \draw[lin] (rCtA) -- (Actuadores);
  \draw[lin] (rSL) -- (Sensores);
  \draw[lin] (rSL) -- (LecturasSensores);
  \draw[lin] (rSE) -- (Sensores);
  \draw[lin] (rSE) -- (EventosInventario);
  \draw[lin] (rAE) -- (Actuadores);
  \draw[lin] (rAE) -- (EventosActuadores);
  \draw[lin] (rEtE) -- (EtiquetasHerramienta);
  \draw[lin] (rEtE) -- (EventosInventario);
  \draw[lin] (rOE) -- (Objeto);
  \draw[lin] (rOE) -- (EtiquetasHerramienta);
  \end{scope}
  \path (rUCo) -- node[card, pos=0.55] {1} (Usuario);
  \path (rUCo) -- node[card, pos=0.55] {N} (ConsultasIA);
  \path (rUE) -- node[card, pos=0.55] {1} (Usuario);
  \path (rUE) -- node[card, pos=0.55] {N} (Equipo);
  \path (rEC) -- node[card, pos=0.55] {1} (Equipo);
  \path (rEC) -- node[card, pos=0.55] {N} (Caja);
  \path (rUC) -- node[card, pos=0.55] {N} (Usuario);
  \path (rUC) -- node[card, pos=0.55] {M} (Caja);
  \path (rUH) -- node[card, pos=0.55] {1} (Usuario);
  \path (rUH) -- node[card, pos=0.55] {N} (Historial);
  \path (rHC) -- node[card, pos=0.55] {1} (Caja);
  \path (rHC) -- node[card, pos=0.55] {N} (Historial);
  \path (rUR) -- node[card, pos=0.55] {1} (Usuario);
  \path (rUR) -- node[card, pos=0.55] {N} (RFID);
  \path (rRA) -- node[card, pos=0.55] {1} (RFID);
  \path (rRA) -- node[card, pos=0.55] {N} (Accesos);
  \path (rCA) -- node[card, pos=0.55] {1} (Caja);
  \path (rCA) -- node[card, pos=0.55] {N} (Accesos);
  \path (rCCt) -- node[card, pos=0.55] {1} (Controladores);
  \path (rCCt) -- node[card, pos=0.55] {1} (Caja);
  \path (rCtS) -- node[card, pos=0.55] {1} (Controladores);
  \path (rCtS) -- node[card, pos=0.55] {N} (Sensores);
  \path (rCtA) -- node[card, pos=0.55] {1} (Controladores);
  \path (rCtA) -- node[card, pos=0.55] {N} (Actuadores);
  \path (rSL) -- node[card, pos=0.55] {1} (Sensores);
  \path (rSL) -- node[card, pos=0.55] {N} (LecturasSensores);
  \path (rSE) -- node[card, pos=0.55] {1} (Sensores);
  \path (rSE) -- node[card, pos=0.55] {N} (EventosInventario);
  \path (rAE) -- node[card, pos=0.55] {1} (Actuadores);
  \path (rAE) -- node[card, pos=0.55] {N} (EventosActuadores);
  \path (rEtE) -- node[card, pos=0.55] {1} (EtiquetasHerramienta);
  \path (rEtE) -- node[card, pos=0.55] {N} (EventosInventario);
  \path (rOE) -- node[card, pos=0.55] {1} (Objeto);
  \path (rOE) -- node[card, pos=0.55] {N} (EtiquetasHerramienta);
\end{tikzpicture}}
\captionof{figure}{Modelo entidad-relación final de la base de datos (simplificado, sin atributos; el modelo completo se muestra en el \ref{anx:er-completo}).}
\label{fig:bd-final}
\end{center}
\begin{multicols}{2}

Haciendo uso de este modelo, se construyó la base de datos en AIVEN, la cual se encuentra en la nube y es accesible desde cualquier lugar. La base de datos contiene todas las tablas necesarias para almacenar la información de los usuarios, equipos, cajas, sensores, actuadores y demás elementos del sistema. Además, se implementaron las relaciones entre las tablas para garantizar la integridad de los datos y facilitar las consultas necesarias para el funcionamiento del sistema. 

Cabe destacar que la base de datos fue transformada para tenerla en la 3FN (Tercera Forma Normal), lo que permite reducir la redundancia de datos y mejorar la eficiencia en las consultas; prevenir anomalías. Esto se logró mediante la normalización de las tablas y la eliminación de dependencias funcionales innecesarias, pasando de un modelo de tablas a la versión que vemos en la versión final.
